% Chapter 2. Literature Review.
% Every statement about a paper was read from the paper itself. Where a number
% is quoted, the paper is cited next to it. Critique paragraphs are marked as
% this project's own reading.

\chapter{Literature Review}
\label{ch:lit}

This chapter reviews the twelve papers on which Phase~1 rests: the model used
(Section~\ref{sec:lit-dai}); the absorber's observed spectrum and height and the mixing
in the clouds, which decide whether a candidate can sit where the absorber is
(Sections~\ref{sec:lit-ph} to~\ref{sec:lit-lef}); and the candidates, with the one paper
that states how strongly the droplets must absorb (Sections~\ref{sec:lit-jiang}
to~\ref{sec:lit-egan}). Section~\ref{sec:lit-gap} draws the threads together. For each
paper, the first part says what the authors did and how, and the second what they found,
followed by a paragraph marked as this project's own reading of what remains unsettled.

Two terms recur. A \emph{mixing ratio} is the fraction of the molecules in a volume of
air that belong to one species, in parts per million (ppmv) or per billion (ppb). The
\emph{eddy diffusion coefficient} $K_{zz}$ is the single number per height with which a
one-dimensional model represents all the vertical stirring it cannot resolve: the larger
$K_{zz}$, the faster a species is mixed up and down. It is quoted in cm$^{2}$\,s$^{-1}$
or m$^{2}$\,s$^{-1}$; one m$^{2}$\,s$^{-1}$ is $10^{4}$~cm$^{2}$\,s$^{-1}$.

% ===========================================================================
\section{Dai et al.\ (2024): the model this project uses}
\label{sec:lit-dai}

\subsection{Objective and Methodology}

Dai, Shao and Sheng built VULCAN-Venus, an open-access one-dimensional photochemistry and
transport model, to study how the chemistry depends on mixing in the clouds and on the
loss of SO\sub{2} into the droplets~\cite{dai2024}. It spans the surface to 112~km in
2~km layers, with 362 neutral and 45 photolysis reactions among 59 species, mainly from
Mills (1998) and Krasnopolsky (2007) (as cited in~\cite{dai2024}). Sunlight is followed
from 100 to 700~nm, and the droplets scatter it in two size modes. The absorber is not a
chemical species: its absorption between 320 and 560~nm is attached to the mode~1
droplets~\cite{dai2024}. The nominal $K_{zz}$, from Bierson and Zhang
(Section~\ref{sec:lit-bz}), rises from about 17\,000~cm$^{2}$\,s$^{-1}$ at the surface to
about 53\,000 at 80~km and steeply above; a case called Model~A replaces it between 46.5
and 55~km with a stronger profile from Venus Express data. Three reactions taken from
Rimmer et al.\ (Section~\ref{sec:lit-rimmer}) move SO\sub{2} into and out of the droplets
as liquid SO\sub{2}, with rates from Henry's law rather than measurement~\cite{dai2024}.

\subsection{Results and Critique}

The model reproduces the observed fall of SO\sub{2} through the clouds: SO\sub{2}
dissolves between 58 and 70~km and is released between 40 and 58~km, where the droplets
evaporate~\cite{dai2024}. Model~A changes the chemistry only slightly. OSSO reaches only
about $10^{-8}$~ppmv at 70~km, which the authors judge insufficient for the absorber; its
photolysis is omitted for lack of laboratory data~\cite{dai2024}.

\emph{This project's reading.} The absorber's amount, height and spectrum are inputs, so
the published model cannot test whether a candidate supplies them; making the absorber
switchable was the first task of this project (Chapter~3). The authors call the liquid
SO\sub{2} a ``substitute medium'' tuned to the observations~\cite{dai2024}, so SO\sub{2}
above the clouds is fitted, not predicted. The mixing test stops at 55~km, below the
cloud top where the absorber acts. And although the abstract speaks of mixing retrieved
from Venus Express, the nominal run uses the Bierson and Zhang profile.

% ===========================================================================
\section{P\'erez-Hoyos et al.\ (2018): the absorber's spectrum and height}
\label{sec:lit-ph}

\subsection{Objective and Methodology}

P\'erez-Hoyos and co-workers analysed 318 spectra from the MASCS spectrometer on
MESSENGER, taken on 5 June 2007 along the equator from near noon to the evening
terminator~\cite{perezhoyos2018}. A radiative transfer retrieval gave the height and
amount of four modes of cloud particles and the absorber's spectral shape. The absorber
was assumed well mixed with mode~1, and only the imaginary part of the mode~1 refractive
index, the part that measures absorption, was retrieved, with SO\sub{2} fixed at
500~ppb. Candidates were then ranked by how closely their spectra follow the result.

\subsection{Results and Critique}

The cloud top, where the optical depth at 0.63~\textmu m reaches one, is at $75 \pm 2$~km;
mode~1 starts at $65 \pm 2$~km with a scale height of about 4~km. The absorption is
centred at $0.34 \pm 0.03$~\textmu m with a full width at half maximum of
$0.14 \pm 0.01$~\textmu m, if Gaussian, and the data are most sensitive to it at
$75 \pm 7$~km~\cite{perezhoyos2018}. In the ranking, S\sub{2}O fits best and OSSO next,
whose shape the authors find closer; FeCl\sub{3}, SCl\sub{2} and Cl\sub{2} absorb over too
narrow a band~\cite{perezhoyos2018}. These values are the observed targets used in this
report (Figure~\ref{fig:F2}).

\emph{This project's reading.} The authors ask a candidate for the right spectrum, a
number density consistent with photochemical models, and a lifetime consistent with the
motion of the ultraviolet markings, but test only the spectrum, as they
state~\cite{perezhoyos2018}. The absorber's height is tied to mode~1 by assumption and
cannot be retrieved separately. The paper also calls the band shifted to longer
wavelengths than an earlier estimate in its results but to shorter ones in its
conclusions; this report uses only the centre and width.

% ===========================================================================
\section{Bierson and Zhang (2020): mixing and the \texorpdfstring{SO\sub{2}}{SO2} problem}
\label{sec:lit-bz}

\subsection{Objective and Methodology}

Bierson and Zhang built the first one-dimensional model to join the lower and middle
atmosphere in one domain, from the surface to 112~km, as a tool for studying how species
and regions interact rather than a tuned fit~\cite{bierson2020}. The absorber is an
imposed opacity, and water and sulfuric acid vapour are fixed because no clouds form.
They compared four $K_{zz}$ profiles: the nominal one, extended downward from an earlier
middle-atmosphere model; one from Krasnopolsky; the nominal divided by ten (Lowkzz); and
the nominal lowered to 1000~cm$^{2}$\,s$^{-1}$ between 45 and 65~km (Cloud
kzz)~\cite{bierson2020}. VULCAN-Venus adopts the nominal profile, as does every Phase~1
run (file \texttt{atm/atm\_Venus\_Bkzz.txt}).

\subsection{Results and Critique}

With the nominal profile SO\sub{2} stays almost uniformly mixed and does not fall above
the clouds~\cite{bierson2020}. Lowkzz and Cloud kzz come close, but the authors call
Lowkzz instructive for understanding the model's sensitivity rather than plausible, and
note that measured and estimated mixing in the clouds
is at least three times stronger than either. They conclude that ``some as of yet unknown
chemical sink in the cloud region or interactions with the cloud droplets'' is
needed~\cite{bierson2020}. OSSO, written (SO)\sub{2}, is below $10^{-25}$ at 58~km, never
exceeds 1~ppb, and peaks near 90~km, far above the cloud top. About 40 per cent of the
rate coefficients have no laboratory value~\cite{bierson2020}.

\emph{This project's reading.} Phase~1 uses a profile whose authors show it cannot
reproduce SO\sub{2} without an extra sink in the clouds, which Dai et al.\ supply. The
profile is published only as a figure, so its values come from the model's input file.
The verdict on OSSO rests on abundance and height, not spectrum.

% ===========================================================================
\section{Rimmer et al.\ (2021): chemistry inside the droplets}
\label{sec:lit-rimmer}

\subsection{Objective and Methodology}

Rimmer and co-workers tested three explanations for the loss of SO\sub{2} and water in
and above the clouds: more water below the clouds, less SO\sub{2} below them, and
chemistry inside the droplets enabled by dissolved hydroxide salts, with sodium hydroxide
as a stand-in~\cite{rimmer2021}. In their 480 species model, gas dissolves into the
droplets by Henry's law and the salts neutralise part of the acid so that more SO\sub{2}
is held. The absorber is a fixed optical depth between 58 and 67~km, fading above.

\subsection{Results and Critique}

Without extra hydrogen, the water available removes only about 20 per cent of the
SO\sub{2}~\cite{rimmer2021}. More water fails, raising the cloud top 20~km too high and
the oxygen far too much. Less SO\sub{2} below the clouds works at 15 to 20~ppm (the paper
gives both), against about 150~ppm observed. Salts reproduce SO\sub{2} with a droplet pH
of about 1 at 50~km, but need a continuous supply of about
$10^{-13}$~mol\,cm$^{-2}$\,s$^{-1}$, some four orders of magnitude more than delivery
from space~\cite{rimmer2021}. No gas in their models explains the absorber above 200~nm,
and the particles' own absorption was not considered.

\emph{This project's reading.} The salt profile is prescribed at each height to fit
SO\sub{2}, its source lies outside the mass balance, and the authors call the solution not
unique~\cite{rimmer2021}. The paper shows that buffered droplets \emph{could} hold the
missing SO\sub{2}, not that a real salt reaches them. Dai et al.\ adopted its droplet
reaction rates.

% ===========================================================================
\section{Lef\`evre et al.\ (2022, 2024): resolved convection in the clouds}
\label{sec:lit-lef}

\subsection{Objective and Methodology}

Lef\`evre and co-workers measured $K_{zz}$ instead of assuming it, with a large-eddy
simulation: a three-dimensional model fine enough to resolve single rising plumes, with a
400~m grid over a box 64~km wide, from the surface to 90~km~\cite{lefevre2022}. The 2022
study followed idealised tracers for SO\sub{2} and water at the equator and at 75 degrees
latitude at noon, and derived $K_{zz}$ from their simulated vertical fluxes. The 2024
study added full photochemistry, equilibrium clouds and midnight, and derived an
effective $K_{zz}$ for each species~\cite{lefevre2024}. The solar heating corresponds to
a fixed, average amount of absorber~\cite{lefevre2022}.

\subsection{Results and Critique}

A deep convective layer sits at about 47 to 56~km at the equator. There, at noon, a cloud-top
convective layer appears at about 66 to 74~km, driven by the absorber's own solar
heating; it is absent at night and at high latitude, and the authors call it still
speculative~\cite{lefevre2022}. The effective $K_{zz}$ in the deep layer is about
$10^{3}$~m$^{2}$\,s$^{-1}$ for SO\sub{2} and $10^{2}$ for water; in the cloud-top layer it
rises from about 1~m$^{2}$\,s$^{-1}$ at 65~km to $10^{4}$ at 72~km, raising SO\sub{2} at
72~km from 0.17 to 0.7~ppm~\cite{lefevre2024}. These values are orders of magnitude above
those in one-dimensional chemistry models~\cite{lefevre2024}, none of which yet includes
the cloud-top convection~\cite{lefevre2022}.

\emph{This project's reading.} Neither paper moves or changes the absorber, though the
2022 paper says its uncertain abundance ``could have a strong effect'' on the cloud-top
convection~\cite{lefevre2022}. Absorber and mixing are coupled: the heating drives mixing
that redistributes the absorber. Neither paper reports mixing above about 75~km, leaving
most of this project's 69 to 93~km response band (Chapter~5) unconstrained, and the
species-dependent $K_{zz}$ shows that VULCAN-Venus's single $K_{zz}$ is a simplification.
These points motivate Phase~2 (Chapter~9).

% ===========================================================================
\section{Jiang et al.\ (2024): iron sulfate minerals}
\label{sec:lit-jiang}

\subsection{Objective and Methodology}

Jiang and co-workers tested in the laboratory which iron sulfate minerals survive in
sulfuric acid of droplet strength, and whether their absorption matches the
absorber~\cite{jiang2024}. Minerals were kept in 50 to 98 weight per cent acid for up to
30 days at room temperature and 1~bar and followed by Raman spectroscopy. The absorbance
of the surviving minerals and of dissolved iron, normalised to a maximum of one between
200 and 500~nm, was fitted to the observed spectrum.

\subsection{Results and Critique}

Rhomboclase is stable up to 74.2~wt\% acid and becomes acid ferric sulfate in stronger
acid: within 14 days at 79.1~wt\%, two hours at 83.5~wt\% and five minutes at
97.4~wt\%~\cite{jiang2024}. The best fit mixes 1~wt\% rhomboclase in 74.2~wt\% acid with
1.25~wt\% acid ferric sulfate in 83.5~wt\% acid; dissolved iron matches 200 to 300~nm.
The authors estimate that about 1 to 2~wt\% of absorber is needed, against at most
$10^{-3}$~wt\% iron from
generous estimates of delivery from space, and that rhomboclase forms only where the
droplets are out of equilibrium or in balance with about 100~ppm or more of water vapour,
above the 30~ppm observed below the clouds~\cite{jiang2024}.

\emph{This project's reading.} Normalised spectra test shape, not strength; the absolute
coefficients are in the supplementary material, which Experiment~13 uses (Chapter~6).
Nothing is modelled about transport, settling or chemistry, so whether the minerals stay
at the cloud top, and how the iron arrives, is untested. Dai et al.\ note that the
minerals formed under a particular laboratory setting, not a Venus-like
one~\cite{dai2024}.

% ===========================================================================
\section{Trabelsi and Francisco (2026): chlorine and sulfur molecules}
\label{sec:lit-trabelsi}

\subsection{Objective and Methodology}

Trabelsi and Francisco calculated from first principles the structure, vibrations,
excited states and cross sections from 200 to 600~nm of three forms of
S\sub{2}Cl\sub{2} (ClSSCl, SSCl\sub{2} and a ring form), as isolated gas molecules at
absolute zero~\cite{trabelsi2026}.

\subsection{Results and Critique}

ClSSCl absorbs strongly near 240~nm, against a measured band near 256 to 260~nm.
SSCl\sub{2} absorbs from about 410~nm, with maxima near 348 and 262~nm, and since this
overlaps the absorber's range the authors conclude it ``could
contribute''~\cite{trabelsi2026}. They offer the results as better inputs for
photochemical models.

\emph{This project's reading.} The paper gives spectra, not amounts: no abundance,
production, lifetime or altitude. A molecule that absorbs in the near ultraviolet is also
broken up by it, so it would need a continuous source. The authors supplied their cross
sections, which Experiment~11 runs in the model (Chapter~6).

% ===========================================================================
\section{Egan et al.\ (2025): OSSO tested end to end}
\label{sec:lit-osso}

\subsection{Objective and Methodology}

Egan and co-workers asked whether OSSO can be abundant enough~\cite{egan2025osso}. They
computed new rate coefficients for its formation from two SO molecules and for its loss,
placed them in a three-dimensional model of the whole atmosphere with a sulfur source from
meteors, ran it for nine Venus days to a steady state, and compared the resulting
reflectance, from a radiative transfer model, with the MESSENGER spectrum.

\subsection{Results and Critique}

OSSO forms a thin layer near 70~km, about three orders of magnitude short of the observed
absorption, still two orders short with its loss to oxygen and chlorine atoms switched
off, and with too narrow a band~\cite{egan2025osso}. SO\sub{2} above the clouds was kept
deliberately too high, which favours OSSO, and the authors conclude that OSSO does not
contribute significantly.

\emph{This project's reading.} This is the only peer-reviewed paper reviewed here in which
the chemistry sets a candidate's amount and the reflectance follows from it, and so the
only independent check on this
project's screen, which reaches the same verdict against a different yardstick
(Experiment~12, Chapter~6). Its radiative step uses one column at the equator near
midday, and its loss rates are calculated, not measured.

% ===========================================================================
\section{Spacek et al.\ (2026): how strongly the droplets must absorb}
\label{sec:lit-spacek}

\subsection{Objective and Methodology}

Spacek, Rimmer, Petkowski and Lee asked how strongly the droplet liquid must absorb, in
units a laboratory measures~\cite{spacek2026}. They fitted a multiple scattering model to
dayside reflectance between 365 and 455~nm from a 2020 campaign, with the absorber spread
evenly through mode~1 and mode~2 droplets between 59 and 65~km and represented by a
lowered \emph{single scattering albedo}, the fraction of light meeting a particle that
is scattered rather than absorbed. The absorption per unit droplet volume gives the
\emph{decadic absorption coefficient} $a$ of the liquid: the absorbance of a 1~cm path
in a laboratory cell.

\subsection{Results and Critique}

The single scattering albedo must be 0.89 at 365~nm, and $a$ peaks at 1278~cm$^{-1}$ at
375~nm, falling below 54~cm$^{-1}$ at 455~nm~\cite{spacek2026}. An efficient organic
absorber would have to make up about 0.5 to 7 per cent of the droplet liquid, nearly four
orders of magnitude above the authors' estimate of the organic carbon available; weaker
inorganic absorbers would have to be a major part of the droplets. Ferric chloride is
rejected because 1~g per litre of \emph{water} gives about 1~cm$^{-1}$ at 375~nm, so a
straight-line extrapolation needs about 1.3~kg per litre~\cite{spacek2026}.

\emph{This project's reading.} The requirement holds only for the assumed layer, mixing
and droplet sizes, as the authors acknowledge. The ferric chloride figure rests on a
measurement in water, though the same paper warns that aqueous spectra have limited value
unless Venus conditions are reproduced~\cite{spacek2026}, and it cites the sulfuric acid
measurements of Egan et al.\ only for spectral shape. Experiment~14, done after Phase~1
closed, repeats the calculation with the same group's measurements of ferric chloride in
hydrochloric acid (Section~\ref{sec:lit-egan} and Chapter~6).

% ===========================================================================
\section{Egan et al.\ (2025, 2026): ferric chloride in sulfuric acid}
\label{sec:lit-egan}

\subsection{Objective and Methodology}

The ferric chloride spectrum used in recent comparisons had been measured in an organic
solvent with which it can react~\cite{egan2025lab}. Egan, James and Plane measured it
from 200 to 600~nm in about 77 and 78~wt\% sulfuric acid with added hydrochloric acid, at
room temperature, and followed its slow conversion to ferric sulfate for about a hundred
days~\cite{egan2025lab}. To fit these mixtures they also derived its molar absorptivity
in hydrochloric acid of several strengths (their Figure~5a), the curves Experiment~14
uses~\cite{egan2025lab}. A preprint by the same group, posted on 14 January 2026 and
marked as not peer reviewed~\cite{egan2026}, added 57 iron reactions to a
three-dimensional model of Venus, with iron from meteors and chlorine from atmospheric
hydrogen chloride, and used the laboratory spectra in a radiative transfer model compared
with the MESSENGER spectrum.

\subsection{Results and Critique}

Ferric chloride had not been measured in concentrated sulfuric acid before, the solvent
has a significant effect on its spectrum, and hydrochloric acid slows its conversion to
sulfate~\cite{egan2025lab}. Its normalised shape agrees well with the absorber, but the
authors add that altitude, formation, lifetime and transport ``must all be considered
through extensive atmospheric modeling'' first. The preprint finds that 1.0 to 1.2~wt\%
ferric chloride in mode~1 droplets reproduces the absorption between 300 and 600~nm (1.2
to 1.5~wt\% if part turns to sulfate), and that a purely meteoric source needs the iron to
stay in the clouds for about 3000 years~\cite{egan2026}.

\emph{This project's reading.} The laboratory spectra were taken several hundred to about
a thousand times more dilute than the fitted loadings, so using them at those loadings
assumes that absorbance grows in proportion to the amount. In the
preprint the amount comes from the radiative fit, not the chemistry, whose run had not
reached a steady state, and the residence time is 2800 and 3400 years in the discussion
but 1800 to 3000 in the conclusions~\cite{egan2026}. Most important here, Spacek et al.\
and the preprint reach opposite conclusions on ferric chloride and neither cites the
other. Experiment~14 (Chapter~6) traces the disagreement to the solvent of the input
measurement; it does not identify the absorber.

% ===========================================================================
\section{What the literature leaves open}
\label{sec:lit-gap}

Against the three tests of Section~\ref{sec:question}, the coverage is uneven.
\textbf{The spectrum} is tested most often, almost always as a normalised shape.
\textbf{The amount} has been tested for a few candidates, one at a time, and only Spacek
et al.\ state the absolute strength the droplets need. \textbf{The altitude} is assumed:
tied to mode~1, set as a fixed layer, or prescribed. No paper reviewed lets a chemistry
model place a candidate by itself and then measures how the radiation responds to where
it sits, while the mixing that would carry a candidate there is uncertain by orders of
magnitude and may be driven, at the cloud top, by the absorber itself. This is the gap of
Section~\ref{sec:objective}, and the following chapters describe how Phase~1 addressed it.
